% Part: proof-theory
% Chapter: cut-elimination
% Section: interpolation

\documentclass[../../../include/open-logic-section]{subfiles}

\begin{document}

\olfileid{pt}{cut}{itp}

\olsection{Maehara's Lemma and Craig's Interpolation Theorem}

The interpolation theorem, due to William Craig, states that if $!A
\Entails !B$ then there is !!a{sentence}~$!C$ (the \emph{interpolant})
such that $!A \Entails !C$ and $!C \Entails !B$. Importantly, the
interpolant~$!C$ can be chosen so that the !!{constant}s and
!!{predicate}s occurring in~$!C$ all occur in both~$!A$ and~$!B$. (We
assume that no !!{function}s are involved.) Without this restriction,
$!A$~and $!B$ could trivially serve as interpolants.

The interpolation theorem holds for classical first-order logic and
has been called the ``last important property of first-order logic''
to be discovered. It has important applications in model theory, and
automated reasoning. Its original motivation and application was in
the philosophy of science, since it can be used to study the
primitives by which a theory can or cannot be axiomatized. It also
implies Beth's definability theorem, which states that if a theory can
implicitly define a concept, it can define it explicitly.

Like many results in mathematical logic, the interpolation theorem can
be proved both using the methods of model theory and those of proof
theory. The advantage of proof theoretic methods is that not only do
they show the existence of an interpolant, but provide an algorithm
that computes the interpolant from !!a{proof} that $!A \Entails !B$,
e.g., !!a{proof} in~\Log{G3c} of $!A \Sequent !B$.

Let $\Lang L(!A)$ be the set of !!{constant}s and !!{predicate}s
in~$!A$ and $\Lang L(\Gamma) = \bigcup \Setabs{\Lang L(!A)}{!A \in
\Gamma}$. Throughout this section, we assume we are dealing with
!!{formula}s without !!{function}s.

\begin{thm}[Craig's Interpolation Theorem]\ollabel{thm:interpolation}
Let the !!{language} of~$!A$ be $\Lang L(!A)$ and that of $!B$ be
$\Lang L(!B)$. If $\Log{G3c} \Proves !A \Sequent !B$ then there is
!!a{formula}~$!C$ in the !!{language}~$\Lang L(!C) \subseteq \Lang
L(!A) \cap \Lang L(!B)$ such that $\Log{G3c} \Proves !A \Sequent !C$
and $\Log{G3c} \Proves !C \Sequent !B$.
\end{thm}

We prove a generalization of the interpolation theorem to make use of
the structure of !!{proof}s in~\Log{G3c}. To this end, imagine the
antecendent and succedents of sequents to be separated into two parts
each. We write $\maeh{\Gamma_1}{\Gamma_2} \Sequent
\maeh{\Delta_1}{\Delta_2}$ to indicate that $\Gamma_1$ and $\Delta_1$
belong to the first part and $\Gamma_2$ and $\Delta_2$ to the second
part. The interpolation theorem then follows immediately from the
following lemma.

\begin{lem}[Maehara's Lemma]\ollabel{lem:maehara} If $\Log{G3c}
\Proves \maeh{\Gamma_1}{\Gamma_2} \Sequent \maeh{\Delta_1}{\Delta_2}$
then there is !!a{formula}~$!C$ in the !!{language}~$\Lang L(!C)
\subseteq \Lang L(\Gamma_1,
\Delta_1) \cap \Lang L(\Gamma_2, \Delta_2)$ such that $\Log{G3c}
\Proves \Gamma_1 \Sequent \Delta_1, !C$ and $\Log{G3c} \Proves !C,
\Gamma_2 \Sequent \Delta_2$.
\end{lem}

\begin{proof}
  Assume that $\pi \Proves \maeh{\Gamma_1}{\Gamma_2} \Sequent
  \maeh{\Delta_1}{\Delta_2}$. We prove the claim by induction on
  $n=\pheight{\pi}$.

  If $n=0$, then $\maeh{\Gamma_1}{\Gamma_2} \Sequent
  \maeh{\Delta_1}{\Delta_2}$ is an axiom, and some atomic
  !!{formula}~$!A$ occurs in both $\Gamma_1, \Gamma_2$ and $\Delta_1,
  \Delta_2$. There are four cases depending on whether $!A$~belongs to
  the first part or the second on the left, and whether it belongs to
  the first or second part on the right.
  \begin{enumerate}
    \item $!A \in \Gamma_1$ and $!A \in \Delta_1$. We want a~$!C$ such
    that both $\Gamma_1 \Sequent \Delta_1, !C$ and $!C, \Gamma_2
    \Sequent \Delta_2$ are !!{provable}. The first is already an axiom
    regardless of how we pick~$!C$ since $!A$ occurs on both the left
    and right. Our only option to guarantee that $!C, \Gamma_2
    \Sequent \Delta_2$ is !!{provable} is to let $!C \ident \lfalse$.
    \item $!A \in \Gamma_1$ and $!A \in \Delta_2$. Then $\Gamma_1
    \Sequent \Delta_1, !A$ and $!A, \Gamma_2 \Sequent \Delta_2$ are
    both axioms, so we can pick $!C \ident !A$.
    \item $!A \in \Gamma_2$ and $!A \in \Delta_1$. Then $!A, \Gamma_1
    \Sequent \Delta_1$ and $\Gamma_2 \Sequent \Delta_2, !A$ would be 
    axioms, but they have~$!A$ on the wrong sides. So $!A$ isn't the
    interpolant. However, using $\RightR{\lnot}$ and $\LeftR{\lnot}$
    we can !!{prove} $\Gamma_1
    \Sequent \Delta_1, \lnot!A$ and $\lnot !A, \Gamma_2 \Sequent
    \Delta_2$.
    \item $!A \in \Gamma_2$ and $!A \in \Delta_2$. Exercise.
  \end{enumerate}
  The cases where $\maeh{\Gamma_1}{\Gamma_2} \Sequent
  \maeh{\Delta_1}{\Delta_2}$ is an axiom because $\lfalse \in
  \Gamma_1$ or $\lfalse \in \Gamma_2$ are left as exercises.

  If $n>0$, $\pi$ ends in a logical inference. We have to distinguish
  cases cases according to which rule was used, and to which part the
  principal formula belongs. In each case we can use the induction
  hypothesis that the lemma holds for all !!{proof}s~$\pi_1$ with
  $\pheight{\pi_1} < n$, and however the end-sequent of~$\pi_1$ may be
  partitioned into two parts. Let's consider some specific cases.
  \begin{enumerate}
    \item $\pi$ ends in $\RightR{\lnot}$ with principal formula~$\lnot
    !A$. We have two cases: depending on whether $\lnot !A$ belongs to
    the first or second part, $\pi$ ends in either
    \[
    \AxiomC{}
    \RightLabel{$\pi_1$}
    \Deduce$\maeh{!A, \Gamma_1}{\Gamma_2} \fCenter
      \maeh{\Delta_1}{\Delta_2}$
    \RightLabel{\RightR{\lnot}}
    \UnaryInf$\maeh{\Gamma_1}{\Gamma_2} \fCenter
      \maeh{\Delta_1, \lnot !A}{\Delta_2}$
    \DisplayProof
    \]
    or
    \[
    \AxiomC{}
    \RightLabel{$\pi_1$}
    \Deduce$\maeh{\Gamma_1}{!A, \Gamma_2} \fCenter
      \maeh{\Delta_1}{\Delta_2}$
    \RightLabel{\RightR{\lnot}}
    \UnaryInf$\maeh{\Gamma_1}{\Gamma_2} \fCenter
      \maeh{\Delta_1}{\Delta_2, \lnot !A}$
    \DisplayProof
    \]
    Note that if the principal !!{formula}~$\lnot !A$ belongs to the
    first part we assigned the active !!{formula}~$!A$ also to the
    first part the premise. This is a choice we get to make. The
    inductive hypothesis would also apply if we assigned the active
    formula to the other part, but then we would not be able to find
    an interpolant (exercise: try it).

    Consider the first situation first, where $\lnot !A$~belongs
    to the first part. By inductive hypothesis, there is an
    interpolant of the end-sequent of~$\pi_1$, i.e., !!a{formula}~$!C_1$
    such that both
    \begin{align*}
    !A, \Gamma_1 & \Sequent \Delta_1, !C_1 &&\text{and} &
    !C_1, \Gamma_2 & \Sequent \Delta_2
    \intertext{are provable. What we need is !!a{formula}~$!C$ such that
    both of the following are !!{provable}:}
    \Gamma_1 & \Sequent \Delta_1, \lnot !A, !C &&\text{and} &
    !C, \Gamma_2 & \Sequent \Delta_2
    \end{align*}
    Clearly, we can just take $!C \ident !C_1$: the inductive hypothesis
    already told us that the sequent on the right is !!{provable}, and
    the sequent on the left follows from the sequent we got from the
    inductive hypothesis by~$\RightR{\lnot}$. The inductive hypothesis
    also tells us that $\Lang L(!C_1) \subseteq \Lang L(!A, \Gamma_1,
    \Delta_1) \cap \Lang L(\Gamma_2, \Delta_2)$. Since $\Lang L(\lnot
    !A) = \Lang L(!A)$ and $\Lang L(!C) = \Lang L(!C_1)$ we have that
    $\Lang L(!C_1) \subseteq \Lang L(\Gamma_1, \Delta_1, \lnot !A) \cap \Lang
    L(\Gamma_2, \Delta_2)$.

    In the second scenario the situation is slightly different
    in that $\lnot !A$~belongs to the second part. We assign the
    active !!{formula} in the premise also to the second part and
    apply the inductive hypothesis to
    \begin{align*}
    \Gamma_1 & \Sequent \Delta_1, !C_1 &&\text{and} &
    !C_1, !A, \Gamma_2 & \Sequent \Delta_2
    \intertext{are provable. If we again apply the \RightR{\lnot} rule to the second of these, we get !!{proof}s of}
    \Gamma_1 & \Sequent \Delta_1, !C_1 &&\text{and} &
    !C_1, \Gamma_2 & \Sequent \Delta_2, \lnot !A
    \end{align*}
    These are again exactly the sequents we want to be !!{provable} if
    $!C_1$ is an interpolant; we again let $!C \ident !C_1$. $\Lang
    L(!C_1) \subseteq \Lang L(\Gamma_1, \Delta_1) \cap \Lang L(\Gamma_2,
    \Delta_2, \lnot !A)$ as before.
    \item $\pi$ ends in $\RightR{\lif}$ with principal !!{formula}~$!A
    \lif !B$. We again have two cases, depending on whether $!A \lif
    !B$ belongs to the first or second part. The !!{proof}~$\pi$ ends
    in either
    \[
    \AxiomC{}
    \RightLabel{$\pi_1$}
    \Deduce$\maeh{!A, \Gamma_1}{\Gamma_2} \fCenter
      \maeh{\Delta_1, !B}{\Delta_2}$
    \RightLabel{\RightR{\lif}}
    \UnaryInf$\maeh{\Gamma_1}{\Gamma_2} \fCenter
      \maeh{\Delta_1, !A \lif !B}{\Delta_2}$
    \DisplayProof
    \]
    or
    \[
    \AxiomC{}
    \RightLabel{$\pi_1$}
    \Deduce$\maeh{\Gamma_1}{!A, \Gamma_2} \fCenter
      \maeh{\Delta_1}{\Delta_2, !B}$
    \RightLabel{\RightR{\lif}}
    \UnaryInf$\maeh{\Gamma_1}{\Gamma_2} \fCenter
      \maeh{\Delta_1}{\Delta_2, !A \lif !B}$
    \DisplayProof
    \]
    Again we have assigned the active !!{formula}s to the same part in
    the premise as the principal !!{formula} in the conclusion. In the
    first secenario, we apply the inductive hypothesis and get
    !!{proof}s of
    \begin{align*}
    !A, \Gamma_1 & \Sequent \Delta_1, !B, !C_1 &&\text{and} & !C_1, \Gamma_2 & \Sequent \Delta_2
    \intertext{The sequent on the left is a suitable premise for \RightR{\lif}. So we have the following !!{provable} sequents:}
    \Gamma_1 & \Sequent \Delta_1, !A \lif !B, !C_1 &&\text{and} & !C_1, \Gamma_2 & \Sequent \Delta_2
    \end{align*}
    This means that $!C_1$ is also an interpolant for the end-sequent
    of~$\pi$ and we can take $!C \ident !C_1$ once more.
    
    The other scenario is handled the same way.
    \item $\pi$ ends in $\LeftR{\lif}$ with principal !!{formula}~$!A
    \lif !B$. If $!A \lif !B$ belongs to the first part, the !!{proof}~$\pi$ ends
    in
    \[
    \AxiomC{}
    \RightLabel{$\pi_1$}
    \Deduce$\maeh{\Gamma_1}{\Gamma_2} \fCenter
      \maeh{\Delta_1, !A}{\Delta_2}$
    \AxiomC{}
    \RightLabel{$\pi_2$}
    \Deduce$\maeh{!B, \Gamma_1}{\Gamma_2} \fCenter
      \maeh{\Delta_1}{\Delta_2}$
    \BinaryInf$\maeh{!A \lif !B, \Gamma_1}{\Gamma_2} \fCenter
      \maeh{\Delta_1}{\Delta_2}$
    \DisplayProof
    \]
    The inductive hypothesis now gives us four !!{provable} sequents,
    two for the interpolant $!C_1$ of the left premise and two for the
    interpolant $!C_2$ of the second premise:
    \begin{align*}
    \Gamma_1 & \Sequent \Delta_1, !A, !C_1 &&\text{(1)} &
    !C_1, \Gamma_2 & \Sequent \Delta_2 && \text{(2)}\\
    !B, \Gamma_1 & \Sequent \Delta_1, !C_2 &&\text{(3)} &
    !C_2, \Gamma_2 & \Sequent \Delta_2 && \text{(4)}
    \end{align*}
    Sequents (1) and (3) can serve as premises of~\LeftR{\lif}, if we
    apply weakening and add $!C_2$ on the right to~(1) and $!C_1$ on
    the right to~(3):
    \[
    \AxiomC{}
    \Deduce$\Gamma_1 \fCenter \Delta_1, !A, !C_1, !C_2$
    \AxiomC{}
    \Deduce$!B, \Gamma_1 \fCenter \Delta_1, !C_1, !C_2$
    \RightLabel{\LeftR{\lif}}
    \BinaryInf$!A \lif !B, \Gamma_1 \fCenter \Delta_1, !C_1, !C_2$
    \DisplayProof
    \]
    By applying $\RightR{\lor}$ we get the sequent
    \[
    !A \lif !B, \Gamma_1 \fCenter \Delta_1, !C_1 \lor !C_2
    \]
    This suggests that $!C_1 \lor !C_2$ might be our interpolant in
    this case. And indeed, we can !!{prove} the other required sequent
    from the remaining sequents (2) and~(4):
    \[
    \AxiomC{}
    \Deduce$!C_1, \Gamma_2 \fCenter \Delta_2$
    \AxiomC{}
    \Deduce$!C_2, \Gamma_2 \fCenter \Delta_2$
    \RightLabel{\LeftR{\lor}}
    \BinaryInf$!C_1 \lor !C_2, \Gamma_2 \fCenter \Delta_2$
    \DisplayProof
    \]
    It remains to verify that $!C_1 \lor !C_2$ is in the correct
    language. The inductive hypothesis gives us that 
    \begin{align*}
    \Lang L(!C_1) & \subseteq
    \Lang L(\Gamma_1, \Delta_1, !A) \cap \Lang L(\Gamma_2, \Delta_2) &\text{and}\\
    \Lang L(!C_2) & \subseteq
    \Lang L(!B, \Gamma_1, \Delta_1) \cap \Lang L(\Gamma_2, \Delta_2)
    \intertext{Since}
    \Lang L(!A) & \subseteq \Lang L(!A \lif !B) &\text{and}\\
    \Lang L(!B) & \subseteq \Lang L(!A \lif !B)
    \intertext{we have both}
    \Lang L(!C_1) & \subseteq
    \Lang L(\Gamma_1, \Delta_1, !A \lif !B) \cap \Lang L(\Gamma_2, \Delta_2) &\text{and}\\
    \Lang L(!C_2) & \subseteq
    \Lang L(\Gamma_1, \Delta_1, !A \lif !B) \cap \Lang L(\Gamma_2, \Delta_2)
    \intertext{and consequently, $\Lang L (!C_1 \lor !C_2) = {}$}
    \Lang L(!C_1) \cup \Lang L(!C_2) & \subseteq
    \Lang L(\Gamma_1, \Delta_1, !A \lif !B) \cap \Lang L(\Gamma_2, \Delta_2),
    \end{align*}
    which is as required.

    In the case where $!A \lif !B$ belongs to the second part, the
    situation is different, but handled in a similar way. The
    inductive hypothesis again gives us four !!{provable} sequents:
    \begin{align*}
    \Gamma_1 & \Sequent \Delta_1, !C_1 &&\text{(1)} &
    !C_1, \Gamma_2 & \Sequent \Delta_2, !A && \text{(2)}\\
    \Gamma_1 & \Sequent \Delta_1, !C_2 &&\text{(3)} &
    !C_2, !B, \Gamma_2 & \Sequent \Delta_2 && \text{(4)}
    \end{align*}
    Now sequents (2) and (4)---plus some weakened !!{formula}s---serve as premises of~\LeftR{\lif}:
    \[
    \AxiomC{}
    \Deduce$!C_1, !C_2, \Gamma_2 \fCenter \Delta_2, !A$
    \AxiomC{}
    \Deduce$!B, !C_1, !C_2, \Gamma_2 \fCenter \Delta_2$
    \RightLabel{\LeftR{\lif}}
    \BinaryInf$!A \lif !B, !C_1, !C_2, \Gamma_2 \fCenter \Delta_2$
    \DisplayProof
    \]
    We again want to turn this into a sequent that has a single
    !!{formula} built from $!C_1$ and $!C_2$, but this time we need it
    in the antecedent (since we're now dealing with the second part).
    Applying $\LeftR{\land}$ does the trick; we should take $!C \ident
    (!C_1 \land !C_2)$ as the interpolant. Fortuitously, we can use
    the remaining sequents (1) and (3) to !!{prove} the corresponding
    sequent for the other part:
    \[
    \AxiomC{}
    \Deduce$\Gamma_1 \fCenter \Delta_1, !C_1$
    \AxiomC{}
    \Deduce$\Gamma_1 \fCenter \Delta_1, !C_2$
    \RightLabel{\RightR{\land}}
    \BinaryInf$\Gamma_1 \fCenter \Delta_1, !C_1 \land !C_2$
    \DisplayProof
    \]
    As before, we have $\Lang L(!C_1 \land !C_2) \subseteq \Lang
    L(\Gamma_1, \Delta_1) \cap \Lang L(\Gamma_2, \Delta_2, !A \lif
    !B)$.

    \item $\pi$ ends in $\RightR{\lforall}$ with principal !!{formula}~$\lforall[x][!A(x)]$:
    \[
    \AxiomC{}
    \RightLabel{$\pi_1$}
    \Deduce$\maeh{\Gamma_1}{\Gamma_2} \fCenter
    \maeh{\Delta_1, !A(c)}{\Delta_2}$
    \RightLabel{\RightR{\lforall}}
    \UnaryInf$\maeh{\Gamma_1}{\Gamma_2} \fCenter
      \maeh{\Delta_1, \lforall[x][!A(x)]}{\Delta_2}$
    \DisplayProof
    \]
    or
    \[
    \AxiomC{}
    \RightLabel{$\pi_1$}
    \Deduce$\maeh{\Gamma_1}{\Gamma_2} \fCenter
      \maeh{\Delta_1}{\Delta_2, !A(c)}$
    \RightLabel{\RightR{\lforall}}
    \UnaryInf$\maeh{\Gamma_1}{\Gamma_2} \fCenter
      \maeh{\Delta_1}{\Delta_2, \lforall[x][!A(x)]}$
    \DisplayProof
    \]
    In the first scenario, the induction hypothesis gives us !!{proof}s
    of $\Gamma_1 \Sequent \Delta_1, !A(c), !C_1$ and $!C_1, \Gamma_2
    \Sequent \Delta_2$. We can apply $\RightR{\lforall}$ to the first
    to obtain $\Gamma_1 \Sequent \Delta_1, \lforall[x][!A(x)], !C_1$.
    (The eigenvariable condition is satisfied, since it was already
    satisfied in the \RightR{\lforall} rule at the end of~$\pi$.) Thus
    $!C_1$ would be an interpolant if the condition that $\Lang
    L(!C_1) \subseteq \Lang L(\Gamma_1, \Delta_1, \lforall[x][!A(x)])
    \cap \Lang L(\Gamma_2, \Delta_2)$ is satisfied. By the inductive
    hypothesis, we have $\Lang L(!C_1) \subseteq \Lang L(\Gamma_1,
    \Delta_1, !A(c)) \cap \Lang L(\Gamma_2, \Delta_2)$. So we have to
    worry about~$c$---can $c$ occur in~$!C_1$? It can't. If $c$ occurs
    in $!C_1$, we would have both $c \in \Lang L(\Gamma_1, \Delta_1,
    !A(c))$ (which is true) and also $c \in \Lang L(\Gamma_2,
    \Delta_2)$. But then $c$ would occur in the conclusion of the
    $\RightR{\lforall}$ inference in~$\pi$, violating the eigenvariable
    condition.
    
    The other scenario is similar.

    \item $\pi$ ends in $\RightR{\lexists}$ with principal
    !!{formula}~$\lexists[x][!A(x)]$. We again have two scenarios. We
    consider the scenario where $\lexists[x][!A(x)]$ belongs to the
    first part, and leave the other as an exercise.

    In the first scenario, $\pi$ ends in
    \[
    \AxiomC{}
    \RightLabel{$\pi_1$}
    \Deduce$\maeh{\Gamma_1}{\Gamma_2} \fCenter
      \maeh{\Delta_1, \lexists[x][!A(x)], !A(t)}{\Delta_2}$
    \RightLabel{\RightR{\lexists}}
    \UnaryInf$\maeh{\Gamma_1}{\Gamma_2} \fCenter
      \maeh{\Delta_1, \lexists[x][!A(x)]}{\Delta_2}$
    \DisplayProof
    \]
    The induction hypothesis gives us !!{proof}s of $\Gamma_1 \Sequent
    \Delta_1, \lexists[x][!A(x)], !A(t), !C_1$ and $!C_1, \Gamma_2
    \Sequent \Delta_2$. We can apply $\RightR{\lexists}$ to the first
    to obtain $\Gamma_1 \Sequent \Delta_1, \lexists[x][!A(x)], !C_1$.
    Thus $!C_1$ would again be an interpolant if the condition that
    $\Lang L(!C_1) \subseteq \Lang L(\Gamma_1, \Delta_1,
    \lexists[x][!A(x)]) \cap \Lang L(\Gamma_2, \Delta_2)$ is
    satisfied. Here the inductive hypothesis gives us $\Lang L(!C_1)
    \subseteq \Lang L(\Gamma_1, \Delta_1, \lexists[x][!A(x)], !A(t))
    \cap \Lang L(\Gamma_2, \Delta_2)$. Now remember that we assumed
    that there are no !!{function}s. That means the term~$t$ can only
    be !!a{constant}, i.e., $t \ident b$. If $b \notin \Lang L(!C_1)$
    or $b \in \Lang L(\Gamma_1, \Delta_1, \lexists[x][!A(x)]) \cap
    \Lang L(\Gamma_2, \Delta_2)$ there is nothing to worry about: we
    can take $!C_1$ as the interpolant. But what if $b \in \Lang
    L(!C_1)$ and $b \notin \Lang L(\Gamma_1, \Delta_1,
    \lexists[x][!A(x)]) \cap \Lang L(\Gamma_2, \Delta_2)$? First of
    all, by replacing all occurrences of~$b$ in~$!C_1$ by
    !!a{variable}~$y$ we can write $!C_1 \ident \Subst{!C_1'}{b}{y}$
    where $!C_1'$~does not contain~$b$. Furthermore, either $b \notin
    \Lang L(\Gamma_1, \Delta_1, \lexists[x][!A(x)])$ or $b \notin
    \Lang L(\Gamma_2, \Delta_2)$. We can take $!C \ident
    \lforall[y][!C_1'(y)]$ or $!C \ident \lexists[y][!C_1'(y)]$,
    respectively. In the first case, we have:
    \[
    \AxiomC{}
    \Deduce$\Gamma_1 \fCenter
      \Delta_1, \lexists[x][!A(x)], !A(b), !C_1'(b)$
    \RightLabel{\RightR{\lexists}}
    \UnaryInf$\Gamma_1 \fCenter \Delta_1, \lexists[x][!A(x)], !C_1'(b)$
    \RightLabel{\RightR{\lforall}}
    \UnaryInf$\Gamma_1 \fCenter
      \Delta_1, \lexists[x][!A(x)], \lforall[y][!C_1'(y)]$
    \DisplayProof
    \]
    and
    \[
    \AxiomC{}
    \Deduce$!C_1'(b), \lforall[y][!C_1'(y)], \Gamma_2 \fCenter \Delta_2$
    \RightLabel{\LeftR{\lforall}}
    \UnaryInf$\lforall[y][!C_1'(y)], \Gamma_2 \fCenter \Delta_2$
    \DisplayProof
    \]
    The !!{proof}s of the top sequents are from the inductive
    hypothesis (in the second one, we apply admissibility of weakening
    with $\lforall[y][!C_1'(y)]$ in the antecedent). The eigenvariable
    condition on \RightR{\lforall} in the first !!{proof} is satisfied
    because $b \notin \Lang L(\Gamma_1, \Delta_1,
    \lexists[x][!A(x)])$, and $!C_1'$~was defined so it does not
    contain~$b$.
    
    In the second case, we have (from the inductive hypothesis plus
    admissibility of weakening by~$\lexists[y][!C_1'(y)]$):
    \[
    \AxiomC{}
    \Deduce$\Gamma_1 \fCenter
      \Delta_1, \lexists[x][!A(x)], !A(b), \lexists[y][!C_1'(y)], !C_1'(b)$
    \RightLabel{\RightR{\lexists}}
    \UnaryInf$\Gamma_1 \fCenter
      \Delta_1, \lexists[x][!A(x)], \lexists[y][!C_1'(y)], !C_1'(b)$
    \RightLabel{\RightR{\lexists}}
    \UnaryInf$\Gamma_1 \fCenter
      \Delta_1, \lexists[x][!A(x)], \lexists[y][!C_1'(y)]$
    \DisplayProof
    \]
    and
    \[
    \AxiomC{}
    \Deduce$!C_1'(b), \Gamma_2 \fCenter \Delta_2$
    \RightLabel{\LeftR{\lexists}}
    \UnaryInf$\lexists[y][!C_1'(y)], \Gamma_2 \fCenter \Delta_2$
    \DisplayProof
    \]
    The eigenvariable condition on \LeftR{\lexists} in the second
    !!{proof} is satisfied because $b \notin \Lang L(\Gamma_2,
    \Delta_2)$, and $!C_1'$ was defined so it does not contain~$b$.
  \end{enumerate}
  The other cases are similary; we leave them as exercises.
\end{proof}

\begin{prob}
Suppose we had a connective $\lnand$ with rules
\[
\Axiom$\Gamma \fCenter \Delta, !A$
\Axiom$\Gamma \fCenter \Delta, !B$
\RightLabel{\LeftR{\lnand}}
\BinaryInf$!A \lnand !B, \Gamma \fCenter \Delta$
\DisplayProof
\quad
\Axiom$!A, !B, \Gamma \fCenter \Delta$
\RightLabel{\RightR{\lnand}}
\UnaryInf$\Gamma \fCenter \Delta, !A \lnand !B$
\DisplayProof
\]
Show that Maehara's Lemma still holds by carrying out the cases in the
proof of \olref{lem:maehara} where $\pi$ ends in one of these rules.
\end{prob}

Having established Maehara's Lemma, we can use it to prove Craig's
Interpolation Theorem.

\begin{proof}[Proof of \olref{thm:interpolation}]
  Suppose $!A \Sequent !B$ is !!{provable}. Apply \olref{lem:maehara}
  to $\maeh{!A}{\ } \Sequent \maeh{\ }{!B}$ to obtain the
  interpolant~$!C$. We have $\Proves !A \Sequent !C$ and $\Proves !C
  \Sequent !B$.
\end{proof}

Suppose $\Gamma$ is a theory (set of !!{sentence}s) in
!!a{language}~$\Lang L$ containing an $n$-place {predicate}~$R$. We
say that $\Gamma$ \emph{implicitly defines}~$R$, iff
\begin{align*}
\Gamma, \Gamma'
& \Entails \lforall[x_1][\dots\lforall[x_n][(R(x_1,\dots,x_n) \liff
R'(x_1, \dots, x_n))]],
\intertext{where $\Gamma'$ is $\Gamma$ with all
occurrences of~$R$ replaced by~$R' \notin \Lang L$. $\Gamma$
\emph{explicitly defines}~$R$ iff there is some~$!A(x_1, \dots, x_n)$
not containing~$R$ such that}
\Gamma &\Entails
\lforall[x_1][\dots\lforall[x_n][(R(x_1,\dots,x_n) \liff !A(x_1,
\dots, x_n))]].
\end{align*}
Clearly, if $\Gamma$ explicitly defines $R$ it implicitly defines~$R$.
The converse is not obvious, but nevertheless holds:

\begin{lem}[Beth Definability Theorem]
  If $\Gamma$ implicitly defines~$R$ then it also explicitly defines it.
\end{lem}

\begin{proof}
  Suppose $\Gamma$ implicitly defines~$R$. Let $R'$, $\Gamma'$ be as
  above and $\Lang L' = \Lang L(\Gamma') = \Lang L \setminus \{R\}
  \cup \{R'\}$. By hypothesis, we have
  \begin{align*}
    \Gamma, \Gamma' & 
    \Sequent \lforall[x_1][\dots\lforall[x_n][(R(x_1,\dots,x_n) 
    \liff R'(x_1, \dots, x_n))]]
  \intertext{Let $c_1$, \dots,~$c_n$ be !!{constant}s not occurring in~$\Gamma$. By the invertibility lemma we have}
    \Gamma, \Gamma', R(c_1,\dots,c_n) & \Sequent  R'(c_1, \dots, c_n)
  \intertext{Apply Maehara's Lemma to}
    \maeh{\Gamma, R(c_1,\dots,c_n)}{\Gamma'} & \Sequent
      \maeh{\ }{R'(c_1, \dots, c_n)}
  \intertext{to find an~$!A \ident !A(x_1, \dots, x_n)$ such that}
    \Gamma, R(c_1,\dots,c_n) & \Sequent !A(c_1, \dots, c_n) \\
    !A(c_1, \dots, c_n), \Gamma' & \Sequent R'(c_1, \dots, c_n)
  \intertext{have !!{proof}s. Since $\Lang L(!A) = \Lang L_1 \cap
  \Lang L_2 = \Lang L(\Gamma) \setminus \{R\}$, $!A$~does not contain~$R$.
  Replace $R'$ everywhere by~$R$ in the !!{proof} of the second sequent
  to get !!a{proof} of}
    !A(c_1, \dots, c_n), \Gamma & \Sequent R(c_1, \dots, c_n)
  \end{align*}
  Apply \RightR{\lif} and \RightR{\lforall} to get !!a{proof} of
  \[
    \Gamma \Sequent \lforall[x_1][\dots\lforall[x_n][(R(x_1,\dots,x_n) \liff !A(x_1, \dots, x_n))]].
  \]
  This shows that $!A(x_1, \dots, x_n)$ explicitly defines $!R$
  in~$\Gamma$.
\end{proof}

A third result which is equivalent to interpolation (and, like it, can
be proved using Maehara's Lemma) is the following theorem: If two
consistent theories~$\Gamma_1$, $\Gamma_2$ contain a complete theory
$\Gamma$ in the language~$\Lang L(\Gamma) \subseteq \Lang L(\Gamma_1)
\cap \Lang L(\Gamma_2)$, then $\Gamma_1 \cup \Gamma_2$ is consistent.

Recall that a complete theory is one such that for every
!!{sentence}~$!A$ in its language either $\Gamma \Proves !A$ or
$\Gamma \Proves \lnot !A$. Since $\Gamma_1$ and~$\Gamma_2$ are consistent
extensions of~$\Gamma$, $\Gamma$ must be consistent. It follows that
if $!A$ is !!a{sentence} in the language $\Lang L(\Gamma_1) \cap \Lang
L(\Gamma_2)$, we can't have $\Gamma_1 \Entails !A$ and $\Gamma_2
\Entails \lnot !A$. To see this, suppose there were such an~$!A$. If
$\Gamma_1 \Entails !A$ then $\Gamma \Entails !A$. Otherwise, by
completeness of~$\Gamma$, $\Gamma \Entails \lnot !A$, and since
$\Gamma \subseteq \Gamma_1$, $\Gamma_1 \Entails \lnot !A$, making
$\Gamma_1$ inconsistent. So $\Gamma \Entails !A$ and since $\Gamma
\subseteq \Gamma_2$, $\Gamma_2 \Entails !A$. So $\Gamma \Entails/
\lnot !A$, or $\Gamma_2$ would be inconsistent.

The non-existence of !!a{sentence}~$!A$ in the joint
!!{language}~$\Lang L(\Gamma_1) \cap \Lang L(\Gamma_2)$ with $\Gamma_1
\Entails !A$ and $\Gamma_2 \Entails \lnot !A$ is all that's needed for
the proof, which we give here proof-theoretically using Maehara's
Lemma.

\begin{thm}[Robinson's Joint Consistency Theorem]
Suppose $\Gamma_1$, $\Gamma_2$ are consistent theories and for no~$!A$
in the !!{language} $\Lang L(\Gamma_1) \cap \Lang L(\Gamma_2)$ do we
have $\Gamma_1 \Entails !A$ and $\Gamma_2 \Entails \lnot !A$. Then
$\Gamma_1 \cup \Gamma_2$ is consistent.
\end{thm}

\begin{proof}
  Suppose for the contrary that $\Gamma_1 \cup \Gamma_2$ is
  inconsistent. Then the sequent $\Gamma_1, \Gamma_2 \Sequent \ $ has
  !!a{proof}. Apply Maehara's Lemma to $\maeh{\Gamma_1}{\Gamma_2}
  \Sequent \maeh{\ }{\ }$ to get !!a{sentence}~$!A$ in the
  !!{language}~$\Lang L(\Gamma_1) \cap \Lang L(\Gamma_2)$ such that
  $\Proves \Gamma_1 \Sequent !A$ and $\Proves !A, \Gamma_2 \Sequent
  \quad$. From the latter we get $\Proves \Gamma_2 \Sequent \lnot !A$.
  But we assumed that no such !!{sentence} exists.
\end{proof}

Above we have tacitly assumed that the theories $\Gamma$, $\Gamma_1$,
$\Gamma_2$ are finite. By compactness, the results also hold for
infinite theories.

\end{document}
